\subsection{Linear Independence in Modules}

Throughout this subsection, $R$ is a nonzero commutative ring with identity $1_R$.\\
Fix an $R$-module $M$ with additive identity $0_M$.

\begin{definition} \label{def:module-linear-independence} \hlt{Linear Independence in a Module}
\begin{enumerate}[label=(\roman*)]
\setlength{\itemsep}{0pt}
\item $S\subseteq M$ is \hlt{$R$-linearly independent} if $(\forall\,F\subseteq S,\ \abs{F}<\infty)(\forall\,(r_v)_{v\in F}\in R^F)\left[\sum\limits_{v\in F}r_vv=0_M\Rightarrow(\forall\,v\in F)\,r_v=0_R\right]$.
\item $S\subseteq M$ is \hlt{$R$-linearly dependent} if it is not $R$-linearly independent.
\end{enumerate}
\end{definition}

\begin{counterexample} \label{cex:module-dependence-need-not-solve-for-vector} \hlt{Dependence Without Solving for One Vector}\\
For distinct $u,v\in M$, dependence of $\{u,v\}$ need not imply $u\in Rv\lor v\in Ru$.\\
In the $\mathbb{Z}$-module $\mathbb{Z}$, the distinct elements $2,3$ satisfy $3\cdot2+(-2)\cdot3=0$.\\
The set $\{2,3\}$ is dependent by Definition \ref{def:module-linear-independence}(i)--(ii) with the nonzero coefficients $3,-2$.\\
Yet $2\notin3\mathbb{Z}$ and $3\notin2\mathbb{Z}$, so neither element is an integer multiple of the other.
\end{counterexample}

\begin{proposition} \label{prop:integral-domain-scaling-preserves-independence} \hlt{Scaling Independent Sets over a Domain}\\
Assume $R$ is integral domain. Let $S\subseteq M$ be finite, $R$-linearly independent,  $r\in R\setminus\{0_R\}$. Then $rS\defeq\{rv:v\in S\}$ is $R$-linearly independent.
\begin{proof}
\begin{enumerate}[label=(\roman*)]
\setlength{\itemsep}{0pt}
\item \hlt{Reduction}: If $S=\emptyset$, then $rS=\emptyset$ is $R$-linearly independent by Definition \ref{def:module-linear-independence}(i).\\
Assume $S\neq\emptyset$ and write $S=\{v_1,\ldots,v_n\}$ with $n\in\N$ and pairwise distinct $v_1,\ldots,v_n$.
\item \hlt{Injective scaling}: If $rv_i=rv_j$ for $i\neq j$, then $rv_i+(-r)v_j=0_M$.\\
Then $r=0_R$ by Definition \ref{def:module-linear-independence}(i), contradicting $r\neq0_R$; hence $rv_1,\ldots,rv_n$ are distinct.
\item \hlt{Relation transfer}: A relation on $T\subseteq rS$ extends to $rS$ by zero coefficients. Fix $(a_1,\ldots,a_n)\in R^n$ satisfying $\sum\limits_{i=1}^{n}a_i(rv_i)=0_M$.\\
$0_M=\sum\limits_{i=1}^{n}a_i(rv_i)=\sum\limits_{i=1}^{n}(ra_i)v_i$ by Definition \ref{def:module-over-ring}(iii), commutativity of $R$. $(\forall\,i\in\{1,\ldots,n\})\,ra_i=0_R$ by Definition \ref{def:module-linear-independence}(i).
\item \hlt{Coefficient cancellation}: $(\forall\,i\in\{1,\ldots,n\})\,a_i=0_R$ by Proposition \ref{prop:integral-domain-cancellation} with $ra_i=0_R=r0_R$ and $r\neq0_R$.\\
Thus $rS$ is $R$-linearly independent by Definition \ref{def:module-linear-independence}(i).
\end{enumerate}
\hlt{Structural Commentary}: Relations on $rS$ reduce to relations on $S$, followed by coefficient cancellation in $R$.\\
\hlt{Key Insight}: Each transferred coefficient has the cancellable factor $r\neq0_R$.
\end{proof}
\end{proposition}
